\section{Topological T-duality and twisted tori}\label{sec:TopTtorus}

In this section we shall apply the results of Section~\ref{sec:NAdef}, and in particular
Green's theorem, to present a scheme
that will be employed in our study of topological T-duality. We shall
then illustrate how our scheme works to reproduce some standard
(commutative) examples of T-dual spaces.

\subsection{Twisted tori and their T-duals}\label{sec:twistedtori}

We are interested in formulating a notion of T-duality for
``torus bundles without $H$-flux'', which for our purposes can be
characterised by the following general class of spaces.

\begin{Definition}[twisted tori]
Let $\sfG$ be a locally compact group which admits a
cocompact discrete subgroup $\Lambda_\sfG$, i.e., a lattice in~$\sfG$, which we let act on $\sfG$ by left
multiplication. The quotient space
\[
\torus_{\Lambda_\sfG}:=\Lambda_\sfG\backslash \sfG
\]
is a \emph{twisted torus}.
\end{Definition}

By~\cite[Lemma~6.2]{Milnor1976}, only unimodular groups can contain lattices, i.e., groups $\sfG$ whose
modular function $\Delta_\sfG$ is identically equal to~$1$.

\begin{Example}[tori]\label{ex:tori}
Every lattice in the abelian Lie group $\sfG=\real^d$ is
isomorphic to $\Lambda_\sfG=\zed^d$, acting by translations. Then
$\torus_{\zed^d}=\zed^d\backslash\real^d=:\torus^d$ is the $d$-dimensional torus.
\end{Example}

\begin{Example}[nilmanifolds]\label{ex:nilmanifolds}
Generalizing Example~\ref{ex:tori}, let $\sfG$ be a connected and simply-connected nilpotent Lie group. Then a theorem of
Malcev~\cite[Theorem~2.12]{Raghunathan1972} establishes the existence
of a lattice $\Lambda_\sfG$ in $\sfG$ if and only if
$\sfG$ can be defined over the rationals,
i.e., there exists a basis for its Lie algebra which has rational
structure constants, and in this case $\torus_{\Lambda_\sfG}=\Lambda_\sfG\backslash \sfG$ is a nilmanifold.
\end{Example}

\begin{Example}[orbifolds]
Let $\sfG$ be a compact Lie group. Then the lattices in $\sfG$ are
precisely the finite subgroups $\Gamma$ of $\sfG$, and
$\torus_\Gamma=\Gamma\backslash\sfG$ is a smooth orbifold. For instance,
for $\sfG=\SU(2)$ the twisted tori are precisely the three-dimensional
ADE orbifolds $\torus_\Gamma=\Gamma\backslash\Sphere^3$
of the three-sphere for a finite subgroup $\Gamma\subset\SU(2)$;
for $\Gamma=\zed_n$ a cyclic subgroup of order $n\geq2$, this recovers
the familiar lens spaces $\torus_{\zed_n}=\zed_n\backslash\Sphere^3=:\mathbb{L}(n,1)$.
\end{Example}

Following~\cite{Mathai2004}, we come now to a central concept of this paper.

\begin{Definition}[topological T-duality]\label{def:Tduality}
Let $\torus_{\Lambda_\sfG}$ be a twisted torus which admits a
non-trivial right action
of the abelian Lie group $\real^n$ for some $n\geq1$. The crossed product
\[
C(\torus_{\Lambda_\sfG})\rtimes_\rt \real^n
\]
is a \emph{$C^*$-algebraic T-dual} of the twisted torus.
\end{Definition}
If the spectrum of the crossed product algebra $
C(\torus_{\Lambda_\sfG})\rtimes_\rt \real^n
$ is a Hausdorff topological space~$X$ (for instance
if it is Morita
equivalent to a commutative $C^*$-algebra $C(X)$), then we say that~$X$ is T-dual to the
twisted torus $\torus_{\Lambda_\sfG}$ and call $X$ a `classical T-dual'; otherwise we say that the T-dual of
$\torus_{\Lambda_\sfG}$ is a noncommutative space.

For Definition \ref{def:Tduality} to be a `good' notion of T-duality, we should first explain
\begin{itemize}\itemsep=0pt
\item[(a)] in
what precise sense $\torus_{\Lambda_\sfG}$ and
$C(\torus_{\Lambda_\sfG})\rtimes_\rt\real^n$ are `equivalent', and
\item[(b)] how
T-duality applied twice returns the original twisted torus.
\end{itemize}
The answers to both of these points turns out to be provided by
working in a suitable category tailored to our treatment of
topological T-duality.

\subsection[T-duality in the category $\CCK\hspace{-1mm}\CCK$]{T-duality in the category $\CCK\hspace{-1mm}\CCK$}
\label{sec:TdualityKK}

The terminology `topological T-duality' refers to a coarse equivalence at
the level of topology; for $C^*$-algebras the topology is measured by
K-theory. A more powerful refinement is provided by Kasparov's bivariant
K-theory which constructs groups ${\rm KK}(\alg,\balg)$ for any pair
of separable $C^*$-algebras $\alg$ and $\balg$; when $\alg=\complex$,
the group ${\rm KK}(\complex,\balg)\simeq {\rm
 K}(\balg)$ is the K-theory group of $\balg$. The cycles in
Kasparov's groups ${\rm KK}(\alg,\balg)$, called \emph{Kasparaov
 $\alg$--$\balg$ bimodules}, are triples $(\hil,\phi,T)$ where $\hil$
is a right Hilbert $\balg$-module, $\phi$ is a $*$-representation of $\alg$
on $\hil$, and $T\in\End_\balg(\hil)$ is a $\balg$-linear operator on
$\hil$, subject to certain compactness conditions; we do not
provide further details of the definition here and instead refer to~\cite{Brodzki2006} for a
concise review of KK-theory in the context that we shall use it
in this paper. Kasparov bimodules may be thought of as generalizations
of morphisms between $C^*$-algebras, in the sense that any algebra homomorphism $\phi\colon \alg\to\balg$
determines a class $[\phi]\in{\rm KK}(\alg,\balg)$, represented by the
$\alg$--$\balg$-bimodule $(\balg,\phi,0)$.

A key feature of Kasparov's KK-theory is the composition product
\[
\otimes_\balg\colon \ {\rm KK}(\alg,\balg)\times{\rm
 KK}(\balg,\Ccal)\longrightarrow{\rm KK}(\alg,\Ccal) ,
\]
which is bilinear and associative. This product is compatible with the
composition of morphisms $\phi\colon \alg\to\balg$ and $\psi\colon \balg\to\Ccal$
of $C^*$-algebras: $[\phi]\otimes_\balg[\psi] = [\psi\circ\phi]$. It
also makes ${\rm KK}(\alg,\alg)$ into a~ring with unit
$1_\alg=[\Id_\alg]$. We say that an element $\alpha\in{\rm
 KK}(\alg,\balg)$ is \emph{invertible} if there exists an element
$\beta\in{\rm KK}(\balg,\alg)$ such that
$\alpha\otimes_\balg\beta=1_\alg$ and
$\beta\otimes_\alg\alpha=1_\balg$.

An important special instance of Kasparov
bimodules comes from Morita equivalence: Any Morita equivalence
$\alg$--$\balg$ bimodule $\Mcal$ is also a Kasparov bimodule
$(\Mcal,\phi,0)$, with $\phi\colon \alg\to\End(\Mcal)$ the left action of
$\alg$, which defines an invertible class $[\Mcal]\in{\rm
 KK}(\alg,\balg)$ with inverse $[\,\overline{\Mcal}\,]\in{\rm
 KK}(\balg,\alg)$ given by the conjugate $\balg$--$\alg$ bimodule
$\overline{\Mcal}$. Generally, if there exists an invertible element
$\alpha\in{\rm KK}(\alg,\balg)$, then the algebras $\alg$ and $\balg$
are said to be \emph{{KK}-equivalent}, and we write $\alg\
\mbf\sim_{\text{\tiny{KK}}}\ \balg$. Thus Morita equivalence implies
KK-equivalence, but the converse is not generally true. KK-equivalent algebras have isomorphic
K-theory groups, but not necessarily homeomorphic spectra.

This refinement naturally
suggests an approach to T-duality where the category of separable
$C^*$-algebras with $*$-homomorphisms is replaced with an additive
category ${\CCK\hspace{-1mm}\CCK}$, whose objects are again separable $C^*$-algebras
but whose morphisms between any two objects $\alg$ and $\balg$ are
given by the classes in ${\rm KK}(\alg,\balg)$ (see,
e.g.,~\cite{Brodzki2007}). The composition product defines the
composition law, and isomorphic algebras
in this category are precisely the KK-equivalent
algebras; in particular, Morita equivalent algebras are isomorphic
as objects in ${\CCK\hspace{-1mm}\CCK}$. Our formulation and computations of topological
T-duality will always take place in this category, and in this
setting we can easily provide answers to points~(a) and~(b) below Definition~\ref{def:Tduality}
through
\begin{Theorem}\label{thm:equivalences}
If $\torus_{\Lambda_\sfG}$ is a twisted torus with a
non-trivial right action
of $\real^n$, then there are
isomorphisms in the category ${\CCK\hspace{-1mm}\CCK}$ given by the equivalences
\begin{itemize}\itemsep=0pt
\item[{\rm (a)}]  $C(\torus_{\Lambda_\sfG}) \
 \mbf\sim_{\text{\tiny{KK}}} \
 C(\torus_{\Lambda_\sfG})\rtimes_\rt\real^n  $ $($up to a shift of degree
 $n \ {\rm mod}~2)$, and
\item[{\rm (b)}]  $\big(C(\torus_{\Lambda_\sfG})\rtimes_\rt
 \real^n\big)\rtimes_{\widehat\rt} \real^n \ \mbf\sim_{\text{\tiny
 M}} \ C(\torus_{\Lambda_\sfG})$.
\end{itemize}
\end{Theorem}
\begin{proof}
The KK-equivalence (a) follows from the Connes--Thom isomorphism,
formulated in the language of KK-theory~\cite{Fack1981}. The Morita
equivalence~(b) follows from Takai duality
(Theorem~\ref{thm:Takaiduality}).
\end{proof}
